\section{Cross-episode memory and governed operator evolution}
\label{sec:fly}
Foundation language models are stateless across sessions: once an interaction ends,
the model retains no memory of whether its previous trading strategy made or lost money.
Fine-tuning the entire LLM with gradient descent after each trading day is impractical,
computationally expensive, and prone to overfitting on noisy financial markets.

To provide cross-episode adaptation without altering LLM weights, we couple an external
associative memory module inspired by the fruit-fly mushroom body \cite{huang2024memory}
with \textbf{Governed Financial Recursive Self-Improvement (\texttt{Fin-RSI})}.
In biological systems, the mushroom body maps complex sensory features into sparse
representations and uses dopamine-gated synaptic plasticity to associate sensory cues
with delayed rewards and punishments across multiple decay rates. Complementarily,
\texttt{Fin-RSI} evolves the mathematical form of the multimodal pooling, subspace
shrinkage, and covariance-adaptation operators across generational research cycles inside
a SHA-256-locked evaluation sandbox.

\paragraph{Associative memory and complete-episode updates.}
Our memory module projects market state features and candidate actions into a high-dimensional
sparse representation (64 Kenyon Cells). When the agent evaluates a new trading decision, it
queries the memory module for associative preference scores based on prior experiences in similar
market regimes. Because financial rewards are delayed over a holding horizon $T$, for positive
account values $W_0,\ldots,W_T$, the target episode return is
\begin{equation}
G = \sum_{t=0}^{T-1}\log(W_{t+1}/W_t)
  = \log(W_T/W_0).
\end{equation}
This return accounts for transaction fees, daily price fluctuations, and final liquidation.
The verification interface checks timestamp ordering, return accounting, and data lineage;
trades relying on look-ahead data or unadjusted prices are rejected before reaching the
memory module, preventing spurious profits from corrupting associative weights. Valid
negative returns remain active learning signals, and each verified trade receipt updates
memory exactly once after settlement clears.

\paragraph{Governed Financial RSI (\texttt{Fin-RSI}) physical operator evolution.}
While associative memory adapts readout weights within a fixed feature parameterization,
raw multimodal financial sequences suffer from three structural pathologies diagnosed by our
physical probes: (i) pre-LayerNorm scalar weight cancellation
($\|\nabla_{w_t}\operatorname{LN}(w_t \mathbf{h}_t)\|_2 = 0$),
(ii) effective sample size ($\operatorname{ESS} = (\sum_t w_t)^2 / \sum_t w_t^2$) collapse during
intraday announcement or forum bursts, and (iii) small-cap subspace variance inflation under
uniform scalar shrinkage. Guided by these probes inside a read-only evaluation harness
(\texttt{HARNESS\_LOCK.json}), \texttt{Fin-RSI} progressively evolves the multimodal operator across
three mathematical generations:
\begin{enumerate}
\item \textbf{Gen-1 (Post-Encoder Value-Space Bounded-Salience \& Burst Conservation):}
To eliminate pre-LayerNorm scale cancellation and bound single-post leverage during intraday
bursts, \texttt{Fin-RSI} moves salience weighting into post-encoder value space with temperature-bounded
activation ($\tau = 2.5$) and Log-Sum-Exp (LSE) intraday burst partition conservation:
\begin{equation}
\begin{aligned}
w_t &= \exp\!\bigl(\tau \tanh(z_t / \tau)\bigr), \\
\bar{\mathbf{v}}_{i,t} &= \frac{\sum_{m \in \mathcal{B}_{i,t}} w_m \mathbf{v}_{i,t,m}}{\sum_{m \in \mathcal{B}_{i,t}} w_m},
\end{aligned}
\end{equation}
lifting sequence effective sample size by $\RSIESSRatio{}\times$ ($\operatorname{ESS}: \RSIUnboundedESS{} \to \RSIBoundedESS{}$).
\item \textbf{Gen-2 (Per-Channel Subspace Precision Gate \& Positive-Part James--Stein Denoising):}
Because fundamental catalyst channels and noisy retail chatter exhibit heterogeneous signal-to-noise
ratios across $D=16$ latent dimensions, \texttt{Fin-RSI} replaces scalar shrinkage with a per-channel
precision gate $\boldsymbol{\alpha}_{i,t,d}$ conditioned on local effective sample size $n_{i,t}^{\text{eff}}$
and cross-modal divergence $\Delta_{i,t,d} = |\mathbf{q}_{i,t,d}^{\text{soc}} - \mathbf{q}_{i,t,d}^{\text{ann}}|$, paired
with positive-part James--Stein empirical-Bayes shrinkage \cite{james1961estimation}:
\begin{equation}
\resizebox{0.86\columnwidth}{!}{$\displaystyle
\begin{aligned}
\alpha_{i,t,d} &= \frac{n_{i,t}^{\text{eff}}}{n_{i,t}^{\text{eff}} + \kappa_d (1 + \Delta_{i,t,d})}, \\
\hat{\theta}_{i,t,d}^{\text{JS}+} &= \mu_{0,d} + \left(1 - \frac{(D-2)\sigma_d^2}{\|\tilde{\mathbf{q}}_{i,t} - \boldsymbol{\mu}_0\|_2^2}\right)_{\!\!+}(\tilde{q}_{i,t,d} - \mu_{0,d}).
\end{aligned}
$}
\end{equation}
\item \textbf{Gen-3 Champion (Streaming Rank-1 Woodbury Covariance Adaptation \& Fisher Gate):}
To capture non-stationary cross-channel correlations without $\mathcal{O}(D^3)$ batch inversion or
look-ahead leakage, \texttt{Fin-RSI} maintains a streaming rank-1 Sherman--Morrison--Woodbury
inverse covariance matrix \cite{sherman1950adjustment} with exponential forgetting factor $\gamma \in (0,1)$
and modulates the 64-KC readout via a Bernoulli Fisher information uncertainty gate
$g_{\text{Fisher}}(p_t) = 4 p_t(1-p_t)$ on the isotonic-calibrated probability $p_t \in (0,1)$:
\begin{equation}
\resizebox{0.86\columnwidth}{!}{$\displaystyle
\begin{aligned}
\boldsymbol{\Sigma}_t^{-1} &= \gamma^{-1}\boldsymbol{\Sigma}_{t-1}^{-1} - \frac{\gamma^{-2}\boldsymbol{\Sigma}_{t-1}^{-1}\mathbf{u}_t\mathbf{u}_t^\top\boldsymbol{\Sigma}_{t-1}^{-1}}{1 + \gamma^{-1}\mathbf{u}_t^\top\boldsymbol{\Sigma}_{t-1}^{-1}\mathbf{u}_t}, \\
g_{\text{Fisher}}(p_t) &= 4 p_t(1-p_t),
\end{aligned}
$}
\end{equation}
where $\mathbf{u}_t = \sqrt{1-\gamma}\,\hat{\boldsymbol{\theta}}_t^{\text{JS}+}$. A candidate generation is
promoted only if it simultaneously passes all 37 counterfactual guards and dominates predecessors across
macro returns, sparse small-cap slices, and crisis regimes.
\end{enumerate}

\paragraph{Isolating learning from risk filtering.}
To verify that performance gains reflect genuine associative learning rather than
passive risk reduction, our benchmark compares the updating memory against a frozen-memory
control and a volatility exposure gate under identical historical splits and costs
(\texttt{FLY\_GATE\_CONTROL\_RESULTS.json}). On a single daily BTC-USDC trajectory ($720$ bars,
$5$ seeds), plastic \texttt{fly} achieves $+5.58\%$ mean test return vs.\ $+2.39\%$ for
\texttt{frozen} ($\Delta = +3.19\%$) and $+3.26\%$ for linear \texttt{ordinary} without the
volatility gate, whereas under the static 80th-percentile volatility gate both \texttt{fly\_gated}
($-0.58\%$) and \texttt{frozen\_gated} ($-0.02\%$) are constrained on the single crypto trend path.
Across $21$ independent market trajectories ($7$ equity boards $\times$ $3$ non-overlapping 3-year
regimes, $2018\text{--}2026$, \PanelBalancedRows{} 2D-balanced rows, $M=5$ seeds) with complete-episode checked
feedback ($G = \log(W_T/W_0)$), holding the 80th-percentile volatility gate, transaction fees
($8\text{ bps}$), and chronological splits strictly identical, \texttt{learn\_with\_gate} achieves
a mean annualized Sharpe of $\mathbf{\FlyPairedLearnSharpe{} \pm 0.14}$ compared to $\mathbf{\FlyPairedFrozenSharpe{} \pm 0.11}$
for \texttt{frozen\_with\_gate} ($\Delta = \mathbf{\FlyPairedDeltaSharpe{} \pm 0.15}$, paired $t = \mathbf{\FlyPairedTStat{}}$,
$p < 10^{-7}$, winning $21/21$ trajectories), $\mathbf{\FlyPairedNoGateSharpe{} \pm 0.16}$ for un-gated \texttt{learn\_no\_gate},
and $\mathbf{\FlyPairedLinearSharpe{} \pm 0.09}$ for dense linear \texttt{ordinary\_with\_gate} (which cannot
separate non-linear contrarian-vs.-momentum regime interactions that sparse 64-KC Kenyon Cell
expansion resolves).

